\addtotoc{Товчлол}
\thispagestyle{plain}

\vspace*{-6mm}
\begin{center}
\textbf{\large \titlename}\\[1mm]
\end{center}

\begin{flushright}
\departmentname\\
Оюутан: \studentID, \authorShort\\
Удирдагч багш: \supervisorname
\end{flushright}

\vspace{1mm}
\setcounter{table}{0}
\setcounter{figure}{0}
\renewcommand{\thetable}{\arabic{table}}
\renewcommand{\thefigure}{\arabic{figure}}
\setlength{\columnsep}{7mm}
\begin{multicols}{2}
\setlength{\parskip}{5pt}
\sloppy

\noindent\textbf{1. Удиртгал}

Энэхүү судалгаа нь олон өнцгийн 2D зургаас 3D загвар үүсгэх бүрэн автоматжуулсан pipeline болон хэрэглэгчийн графикийн интерфейсийг (GUI) боловсруулах, мөн matching параметрүүдийн реконструкцийн нарийвчлалд үзүүлэх нөлөөг туршилтаар тодорхойлох асуудлыг авч үзсэн. Уламжлалт Structure-from-Motion (SfM) болон Multi-View Stereo (MVS) аргууд нь геометрийн нарийвчлалаараа давуу боловч тооцооллын шаардлага өндөр, тохиргооны сонголт нь гаралтын чанарт шийдвэрлэх нөлөөтэй байдаг.

Судалгааны зорилго нь: COLMAP болон OpenMVS-д суурилсан 9 үе шаттай pipeline боловсруулах; PyQt5-д суурилсан GUI систем бүтээх; matcher болон хослолын аргын нарийвчлалд үзүүлэх нөлөөг туршилтаар үнэлэх; Chamfer Distance болон RMSE-ээр реконструкцийн чанарыг хэмжих; судалгааны хязгаарлалтыг ил тод тэмдэглэх. Үнэлэлтэд Chamfer Distance ба RMSE ашигласан бөгөөд утга бага байх тусам сэргээгдсэн загвар ground truth-тай илүү ойр байна.

\noindent\textbf{2. Судалгааны ажлын онол, өнөөгийн түвшин}

Гол урьдчилсан судалгааны ажлуудаас: Hartley \& Zisserman (2003) нар олон камерын геометрийн суурь онолыг тавьсан; Lowe (2004)-ийн SIFT нь масштаб болон эргэлтэд тэсвэртэй фичер илрүүлэлтийг нэвтрүүлсэн; Sch{\"o}nberger \& Frahm (2016)-ийн COLMAP нь incremental SfM pipeline болсон; Cernea (2015)-ийн OpenMVS нь нягт реконструкцийн нэгдсэн хэрэгсэл болсон. NeRF (Mildenhall et al., 2020) болон 3D Gaussian Splatting (Kerbl et al., 2023) нь дүрслэлийн чанараараа давуу боловч их хэмжээний тооцооллын нөөц шаарддаг. Энэхүү судалгаанд найдвартай байдал болон ирээдүйд машин сургалтын аргуудыг нэгтгэх уян хатан байдлыг харгалзан уламжлалт COLMAP+OpenMVS аргыг сонгосон.

\noindent\textbf{3. Судалгааны арга зүй ба хэрэгжүүлэлт}

Pipeline нь Python програмчлалын хэл дээр (\texttt{core/}, \texttt{pipeline/}, \texttt{gui/}, \texttt{utils/} модуль) бичигдсэн. PyQt5-д суурилсан GUI нь дэвшлийн мөр, бодит цагийн лог болон pipeline-ийг зогсоох боломжийг хэрэглэгчид олгодог. Техник хангамж: NVIDIA RTX 4060, Intel Core i9 (15-р үе), 16\,GB RAM.

Туршилтад гурван датасет ашигласан (Хүснэгт~\ref{tab:tovchlol_datasets}). Temple Ring болон Caterpillar датасетууд нь $1920\times1080$ нягтралтай JPEG форматтай объектын зургуудаас бүрдэнэ. 3-р туршилтад гадаа орчны бодит дүрс сэргээлтийн туршилт болгон нийтэд нээлттэй 4K Lighthouse видео (5:45, 23.98\,fps, $4096\times2160$) ашигласан бөгөөд уг бичлэгийг зохиогчид бүтээгээгүй.

\begin{table}[H]
\centering
\caption{Туршилтын датасетууд}
\label{tab:tovchlol_datasets}
\scriptsize
\setlength{\tabcolsep}{3pt}
\begin{tabular}{|l|r|l|}
\hline
\textbf{Датасет} & \textbf{Зураг} & \textbf{Ашигласан} \\ \hline
Temple Ring & 47 & Туршилт 1 \\ \hline
Caterpillar & 380${\sim}$ & Туршилт 2 \\ \hline
Lighthouse 4K видео & ${\sim}$86/172/200 фрейм & Туршилт 3 \\ \hline
\end{tabular}
\end{table}

Гурван туршилтын цуврал явуулсан:
\begin{itemize}
    \item \textbf{1-р туршилт:} SIFT болон LightGlue matcher-ийг exhaustive болон sequential (overlap=10) хослолын аргатай хослуулан 4 тохиргоогоор Dataset A-д харьцуулсан.
    \item \textbf{2-р туршилт:} Sequential matching-ийн overlap параметрийг 5, 10, 20, 40 гэсэн утгаар Dataset A-д туршиж оновчтой босго тодорхойлсон.
    \item \textbf{3-р туршилт:} Lighthouse видеогоос 4с/2с/1.5с давтамжаар (${\sim}86$, ${\sim}172$, ${\sim}200$ фрейм) гаргаж авсан.
\end{itemize}

\noindent\textbf{4. Судалгааны үр дүн, нэгтгэсэн дүгнэлт}

Системийн гол гүйцэтгэлийн үзүүлэлтүүд:

\begin{table}[H]
\centering
\caption{Системийн гол гүйцэтгэлийн үзүүлэлтүүд}
\label{tab:tovchlol_results}
\scriptsize
\setlength{\tabcolsep}{3pt}
\begin{tabular}{|l|r|}
\hline
\textbf{Үзүүлэлт} & \textbf{Утга} \\ \hline
Temple Ring (зураг) & 47 \\ \hline
Sparse point cloud & 35{,}000+ цэг \\ \hline
Dense point cloud & 2.1 сая цэг \\ \hline
Mesh vertices / faces & 1.2 сая / 2.3 сая \\ \hline
Хамгийн сайн CD (1B: LG+Exh.) & 0.0074\,м \\ \hline
Хамгийн сайн RMSE (1A: SIFT+Exh.) & 0.0064\,м \\ \hline
\end{tabular}
\end{table}

\textbf{1-р туршилт} нь matching стратегийн сонголт (exhaustive эсвэл sequential) нь matcher-ийн сонголтоос илүү чухал нөлөөтэй болохыг харуулсан. 1A болон 1B бараг ижил үр дүн үзүүлсэн бол 1D (LightGlue+Sequential) гамшигт бүтэлгүйтсэн --- overlap=10 sequential matching нь 47 зургийн датасетэд LightGlue-д хангалттай тохирол олж чадаагүй.
\end{multicols}

\begin{figure}[H]
\centering
\begin{tikzpicture}
\begin{axis}[
    xbar, bar width=12pt,
    width=0.82\textwidth, height=5cm,
    symbolic y coords={1A,1B,1C,1D},
    ytick=data,
    yticklabels={
        {1A: SIFT+Exhaustive},
        {1B: LG+Exhaustive},
        {1C: SIFT+Sequential},
        {1D: LG+Sequential}},
    xlabel={Chamfer Distance (м)},
    xmin=0, xmax=0.17,
    enlarge y limits=0.3,
    grid=major, grid style={dashed,gray!30},
]
\addplot[fill=blue!60,draw=blue!80] coordinates {
    (0.007803,1A)(0.007437,1B)(0.010782,1C)(0.143084,1D)
};
\end{axis}
\end{tikzpicture}
\caption{1-р туршилт --- matcher болон хослолын аргын Chamfer Distance харьцуулалт}
\label{fig:tovchlol_exp1_cd}
\end{figure}

\begin{figure}[H]
\centering
\begin{tikzpicture}
\begin{axis}[
    xbar, bar width=12pt,
    width=0.82\textwidth, height=5cm,
    symbolic y coords={1A,1B,1C,1D},
    ytick=data,
    yticklabels={
        {1A: SIFT+Exhaustive},
        {1B: LG+Exhaustive},
        {1C: SIFT+Sequential},
        {1D: LG+Sequential}},
    xlabel={RMSE (м)},
    xmin=0, xmax=0.31,
    enlarge y limits=0.3,
    grid=major, grid style={dashed,gray!30},
]
\addplot[fill=red!55,draw=red!75] coordinates {
    (0.006362,1A)(0.007368,1B)(0.009250,1C)(0.278657,1D)
};
\end{axis}
\end{tikzpicture}
\caption{1-р туршилт --- matcher болон хослолын аргын RMSE харьцуулалт}
\label{fig:tovchlol_exp1_rmse}
\end{figure}

\begin{multicols}{2}

\textbf{2-р туршилт} нь sequential overlap-д тодорхой босгын нөлөө байгааг тогтоосон. overlap=5-аас overlap=10 руу шилжихэд Chamfer Distance 0.033\,м-ээр буурсан; 10-аас дээш нэмэгдүүлэх нь бараг ямар ч ахиц өгөхгүй байв. Overlap=10 нь энэ төрлийн датасетэд оновчтой тохиргоо болохыг баталгаажуулсан.

\end{multicols}

\begin{figure}[H]
\centering
\begin{tikzpicture}
\begin{axis}[
    ybar, bar width=18pt,
    width=0.6\textwidth, height=5.5cm,
    symbolic x coords={5,10,20,40},
    xtick=data,
    xlabel={Sequential overlap утга},
    ylabel={Алдаа (м)},
    ymin=0, ymax=0.18,
    enlarge x limits=0.2,
    legend style={at={(0.97,0.97)},anchor=north east,font=\small},
    grid=major, grid style={dashed,gray!30},
]
\addplot[fill=blue!60,draw=blue!80] coordinates {
    (5,0.152750)(10,0.119549)(20,0.120356)(40,0.119830)
};
\addplot[fill=red!55,draw=red!75] coordinates {
    (5,0.123995)(10,0.078579)(20,0.084193)(40,0.080573)
};
\legend{Chamfer Distance, RMSE}
\end{axis}
\end{tikzpicture}
\caption{2-р туршилт --- Sequential overlap болон алдааны харьцаа. overlap=5-аас 10 руу шилжихэд чанар эрс сайжирч, 10-аас дээш ахиц бага.}
\label{fig:tovchlol_exp2}
\end{figure}

\begin{multicols}{2}

\textbf{3-р туршилт}эд нийтэд нээлттэй Lighthouse 4K видеог гадаа орчны оролт болгон ашигласан. Гол арга зүйн олдвор нь 3A-ийн бага Chamfer Distance нь сайн реконструкцийн чанарыг заагаагүй --- эталон цэгэн үүлний хэмжээний 18 дахин ялгаа (952к ба 17.5М цэг) нь тохиргоо хоорондын шууд харьцуулалтыг найдваргүй болгосон. Харьцуулах боломжтой 3B болон 3C-ийн хооронд фрейм нягтралыг нэмэгдүүлэх нь Chamfer Distance-ийг 0.012\,м-ээр бууруулсан.

\end{multicols}

\begin{figure}[H]
\centering
\begin{tikzpicture}
\begin{axis}[
    ybar, bar width=20pt,
    width=0.6\textwidth, height=5.5cm,
    symbolic x coords={3A (${\sim}$86),3B (${\sim}$172),3C (${\sim}$200)},
    xtick=data,
    xlabel={Тохиргоо (фреймийн тоо)},
    ylabel={Алдаа (м)},
    ymin=0, ymax=0.095,
    enlarge x limits=0.2,
    legend style={at={(0.97,0.97)},anchor=north east,font=\small},
    grid=major, grid style={dashed,gray!30},
]
\addplot[fill=blue!60,draw=blue!80] coordinates {
    (3A (${\sim}$86),0.010333)
    (3B (${\sim}$172),0.053916)
    (3C (${\sim}$200),0.041948)
};
\addplot[fill=red!55,draw=red!75] coordinates {
    (3A (${\sim}$86),0.017397)
    (3B (${\sim}$172),0.082545)
    (3C (${\sim}$200),0.066822)
};
\legend{Chamfer Distance, RMSE}
\end{axis}
\end{tikzpicture}
\caption{3-р туршилт --- Фрейм сонголт болон алдаа (Lighthouse 4K видео). Эталон цэгэн үүлний 18 дахин ялгаанаас болж 3A болон 3B/3C шууд харьцуулахад тохиромжгүй.}
\label{fig:tovchlol_exp3_errors}
\end{figure}

\begin{figure}[H]
\centering
\begin{tikzpicture}
\begin{axis}[
    ybar, bar width=20pt,
    width=0.6\textwidth, height=4.5cm,
    symbolic x coords={3A (${\sim}$86),3B (${\sim}$172),3C (${\sim}$200)},
    xtick=data,
    xlabel={Тохиргоо (фреймийн тоо)},
    ylabel={Эталон цэг (сая)},
    ymin=0, ymax=24,
    enlarge x limits=0.2,
    grid=major, grid style={dashed,gray!30},
]
\addplot[fill=teal!55,draw=teal!75] coordinates {
    (3A (${\sim}$86),0.952)
    (3B (${\sim}$172),17.468)
    (3C (${\sim}$200),21.452)
};
\end{axis}
\end{tikzpicture}
\caption{3-р туршилт --- Эталон цэгэн үүлний хэмжээ. 3A болон 3B/3C хоорондын 18 дахин ялгаа нь шууд метрийн харьцуулалтыг хүчингүй болгодог.}
\label{fig:tovchlol_exp3_pts}
\end{figure}

\begin{multicols}{2}

Дараах хязгаарлалтыг тэмдэглэсэн:
\begin{itemize}
    \item Зөвхөн нэг объектын төрлийг туршиж үзсэн тул (1 ба 2-р туршилт) олон объектод ажиллах чадварыг баталгаажуулаагүй,
    \item 3-р туршилтын эталон цэгэн үүлний жигд бус хэмжээ нь тохиргоо хоорондын харьцуулалтыг найдваргүй болгодог,
    \item GUI-д камерын калибрацийн интерфейс дутмаг,
    \item Sequential matching нь жижиг датасетэд зарим matcher-т бүтэлгүйтдэг.
\end{itemize}

Инженерчлэлийн хувьд энэ судалгаа нь бүх реконструкцийн pipeline-ийг нэг цогц PyQt5 GUI системд нэгтгэсэн. Цаашдын судалгаанд SuperPoint/SuperGlue болон LightGlue-г нэгтгэх, NeRF болон Gaussian Splatting-тэй харьцуулах, олон янзын объектын ангилалд үнэлэлтийг өргөтгөх нь чухал алхам болно.

\noindent\textbf{5. Ном зүй}
\small
\begin{enumerate}
\setlength{\itemsep}{1pt}\setlength{\parsep}{0pt}\setlength{\topsep}{2pt}
    \item Hartley, R., \& Zisserman, A. (2003). \textit{Multiple View Geometry in Computer Vision}. Cambridge University Press.
    \item Lowe, D. G. (2004). Distinctive image features from scale-invariant keypoints. \textit{IJCV}, 60(2), 91--110.
    \item Sch{\"o}nberger, J. L., \& Frahm, J. M. (2016). Structure-from-motion revisited. \textit{CVPR}, 4104--4113.
    \item Cernea, D. (2015). OpenMVS: Multi-View Stereo Reconstruction Library.
    \item Mildenhall, B., et al. (2020). NeRF: Representing scenes as neural radiance fields. \textit{ECCV}, 405--421.
    \item Lindenberger, P., et al. (2023). LightGlue: Local feature matching at light speed. \textit{ICCV}, 17627--17638.
\end{enumerate}
\normalsize

\end{multicols}
